\documentclass[10pt,letterpaper,sans]{moderncv}
\moderncvstyle{banking}
\moderncvcolor{blue}

\usepackage[scale=0.84]{geometry}
\nopagenumbers{}
\setlength{\hintscolumnwidth}{2.35cm}

\name{Krishnaa}{Vadivel}
\title{Electrical Engineer | Semiconductor Devices | Photonics | Embedded Systems}
\email{srikrishnaa.careers@gmail.com}
\phone[mobile]{516-853-5663}
\social[linkedin]{sri-krishnaa-ece-phd}
\extrainfo{\href{https://krishnaa423.github.io/personal-website/}{Website} \textbar{} \href{https://github.com/krishnaa423}{GitHub: krishnaa423} \textbar{} \href{https://leetcode.com/u/krishnaa42342}{LeetCode: krishnaa42342} \textbar{} \href{https://hub.docker.com/u/krishnaa42342}{Docker Hub: krishnaa42342}}

\begin{document}
\makecvtitle

\section{Profile}
\cvitem{}{Electrical engineering PhD researcher with broad preparation across thin-film fabrication, semiconductor physics, integrated photonics, device simulation, and scientific software. Combines cleanroom and characterization experience with light-matter modeling, distributed HPC software, and machine learning for semiconductor and photonic systems.}

\section{Research Experience}
\cventry{2021--present}{Research Assistant / PhD Researcher}{Yale University, Electrical Engineering}{}{}{\begin{itemize}%
\item Developed GaN-based power-electronics devices in cleanroom and MOCVD environments during the early PhD, including fabrication-oriented process work, selective-area etching techniques, and materials development.
\item Characterized semiconductor materials with SEM, cathodoluminescence, and Raman spectroscopy, connecting fabrication outcomes to optical and structural behavior.
\item Built simulation tools for VCSEL emission and coupled photon, phonon, and electron interactions in semiconductor and photonics problems.
\item Developed distributed HPC software for exciton-polaron and light-matter interaction modeling using first-principles methods, quantum many-body theory, MPI, OpenMP, CUDA, ROCm, PETSc, and SLEPc.
\item Used agentic ML and retrieval-based tooling to generate first-principles software inputs from journal literature and technical documentation.
\item Maintained strong continuity between theory and engineering practice through coursework and self-developed notes in electromagnetism, signals/systems, semiconductor transport, optics, and device equations.
\item Co-authored a 2025 \textit{Journal of the American Chemical Society} paper and presented related APS work in 2024, 2025, and 2026 on self-trapped excitons and excited-state materials phenomena.
\end{itemize}}

\section{Professional Experience}
\cventry{2019--2021}{Software and Embedded Systems Engineer}{Probe Technologies Inc. / Weatherford Wireline Services}{}{}{\begin{itemize}%
\item Parallelized CUDA-based software for analysis and visualization of acoustic and nuclear well-tool data used in engineering diagnostics and field interpretation.
\item Applied dolfinx-based finite-element analysis to frequency studies of steel pipes in well environments to support physics-driven engineering analysis.
\item Contributed instrumentation-oriented engineering around deployed well tools, connecting field data, sensing systems, and numerical analysis.
\end{itemize}}
\cventry{2018--2019}{Photonics Technical Writer}{FindLight}{}{}{\begin{itemize}%
\item Produced technical articles on lasers, photonics components, optical systems, and engineering applications for a specialized industrial audience.
\end{itemize}}

\section{Selected Projects}
\cventry{}{Integrated Photonics Band-Pass Filter}{}{}{}{\begin{itemize}%
\item Designed a Meep-based integrated photonics band-pass filter, then fabricated and tested the device with an erbium-doped fiber laser setup.
\end{itemize}}
\cventry{}{Semiconductor Device Simulation}{}{}{}{\begin{itemize}%
\item Developed numerical treatments of Poisson, drift-diffusion, and pn-junction-style device equations to connect semiconductor theory with engineering-oriented simulation.
\end{itemize}}
\cventry{}{VCSEL Light-Matter Modeling}{}{}{}{\begin{itemize}%
\item Built simulation tools for VCSEL emission and coupled photon, phonon, and electron interactions in semiconductor systems.
\end{itemize}}
\cventry{}{Equivariant Neural Networks for Materials and Optics}{}{}{}{\begin{itemize}%
\item Developed PyTorch Geometric equivariant neural networks for material energy prediction and optical-property prediction.
\end{itemize}}

\section{Education}
\cventry{2021--present}{PhD, Electrical Engineering}{Yale University}{}{}{Ongoing research spanning computational materials, semiconductor-relevant modeling, and scientific software.}
\cventry{2023}{MPhil, Electrical Engineering}{Yale University}{}{}{}
\cventry{2023}{MS, Electrical Engineering}{Yale University}{}{}{}
\cventry{2017--2018}{MEng, Electrical and Computer Engineering}{Cornell University}{}{}{Included work in stochastic computing, semiconductor/device topics, and advanced EE design.}
\cventry{2013--2017}{BS, Electrical and Computer Engineering}{Cornell University}{}{}{Included strong preparation in embedded systems, robotics, controls, and hardware-software integration.}

\section{Selected Coursework}
\cvitem{Semiconductors}{Semiconductor Physics; Thin Film Science; drift-diffusion and device modeling foundations}
\cvitem{Photonics}{Lasers and Optoelectronics; Integrated Photonics}
\cvitem{Circuits and VLSI}{VLSI Design; Analog VLSI Design; stochastic-computing-related design work}
\cvitem{Systems}{Mathematics of Signals and Systems; Advanced Embedded Systems}
\cvitem{Foundations}{Graduate Quantum Mechanics; Solid State Physics; Many-Body Quantum Physics}

\section{Technical Skills}
\cvitem{Languages}{Python, C, C++, CUDA, JavaScript, Bash, PowerShell}
\cvitem{Fabrication / Characterization}{Cleanroom processing, MOCVD, SEM, cathodoluminescence, Raman spectroscopy, laser-bench optical testing}
\cvitem{Simulation}{Semiconductor device modeling, FDTD photonics simulation, finite-element analysis, first-principles modeling, quantum many-body modeling}
\cvitem{Machine Learning}{PyTorch, PyTorch Geometric, equivariant graph neural networks, materials- and optical-property prediction, agentic input generation}
\cvitem{Scientific Computing}{MPI, OpenMP, CUDA, ROCm, PETSc, SLEPc, NumPy, pandas, SciPy, SymPy, matplotlib}
\cvitem{Software / Platforms}{Python, C, C++, C\#, ASP.NET, Microsoft SQL Server, Linux command line, Docker, Docker Compose, Kubernetes, gcloud, Frontier, Frontera, Perlmutter}

\end{document}
